\section{Main Theorem}
\label{sec:main-theorem}

\begin{theorem}[Conditional canonical-catalog representation]
\label{thm:main}
Under Assumptions~\ref{assump:sq-parameter-regime},
\ref{assump:universal-adversarial-sq}, and
\ref{assump:canonical-rounded-output-catalog}, the one deterministic map
\(\phi:\mathcal X\to\mathbb R^L\) in \eqref{eq:main-notation} is fixed
independently of the current distribution, target, valid oracle policy, and
learner tape, and
\begin{equation}
\forall h\in\mathcal H\ \exists w_h\in\Delta_L\ \forall x\in\mathcal X,
\qquad
h(x)\langle w_h,\phi(x)\rangle
\ge1-2\varepsilon>\frac12>0.
\label{eq:main-margin}
\end{equation}
Only \(h\) may affect \(w_h\).  Consequently,
\begin{equation}
\operatorname{dc}(\mathcal H)
\le L
\le B\left(1+\frac{m}{\tau^2}\right)^k.
\label{eq:main-dimension-bound}
\end{equation}
If \(\mathcal H=\varnothing\) or \(\mathcal X=\varnothing\), then in fact
\(\operatorname{dc}(\mathcal H)=0\).

The conclusion is deterministic and pointwise, with no exceptional set.  It
uses the fixed finite horizon \(m\), including \(m=0\), and has no hidden
constant: all dependence on \(m,\tau,\varepsilon,L,B,k\) is displayed, while
\(B,k\) are the fixed protocol-family constants described in
Assumption~\ref{assump:sq-parameter-regime}.  The result is conditional on
the canonical-policy-only catalog in
Assumption~\ref{assump:canonical-rounded-output-catalog}; it asserts no
catalog property for other valid policies.
\end{theorem}
